\section{How to use this document}
\label{sec:how-to-use}

\subsection{Purpose}
\begin{itemize}
  \item Survive context resets (new chats, compacted histories).
  \item Record \textbf{why}, not only \textbf{what}.
  \item Describe the \textbf{current} architecture (token renderer, CA sim, tests)
        and any remaining polish work.
  \item Separate \emph{observed reality} from \emph{provisional} notes.
\end{itemize}

\subsection{Conventions}
\begin{description}[style=nextline]
  \item[Status boxes]
    Snapshot of maturity (prototype, shipped, production-ish).
  \item[Decision boxes]
    Closed choices. Prefer adding a new decision over rewriting history silently;
    note supersession dates.
  \item[Open boxes]
    Questions only the author can settle.
  \item[Inference boxes]
    Reverse-engineered claims. Treat as provisional.
  \item[\texttt{\textbackslash todoauthor\{...\}}]
    Explicit author TODO inline in the prose.
\end{description}

\subsection{Workflow for future sessions}
\begin{enumerate}
  \item Read the repository and closest nested \texttt{AGENTS.md} files for
        operational rules.
  \item Read \pathref{architecture-consensus.md} under \pathref{docs/}, the \textbf{single
        unified} architecture document (purpose, historical plans map, code
        truth, decisions, bugs, debt, work order).
  \item Optionally skim \cref{sec:overview}, \cref{sec:rendering-current},
        \cref{sec:rendering-target}, \cref{sec:game}, and the decision log
        (\cref{sec:decisions}) for long-form detail.
  \item Use the records under \pathref{docs/sessions/} for dated implementation
        and verification evidence. Treat \pathref{docs/history/} as provenance.
  \item Implement against those sources, not against chat memory or old PDFs alone.
  \item When a decision is made, append a dated entry to the decision log
        \textbf{and} update the consensus Markdown.
  \item When code lands, update the consensus file and the ``current state''
        sections. If docs and code disagree, \textbf{code wins} until docs catch up.
\end{enumerate}

\subsection{Building the PDF}
From the repo root (requires PowerShell and a TeX distribution with
\texttt{latexmk}):
\begin{lstlisting}[language={}]
.\docs\build.ps1
\end{lstlisting}

\noindent Outputs: \pathref{docs/output/illumo.pdf} and
\pathref{docs/output/architecture-map.pdf}.

\begin{statusbox}
As of 2026-08-19 (current-state audit): engine, scene, system, and platform code is
under \pathref{Illumo/Include} and \pathref{Illumo/Source}; IllumoGame contains
only game and ruleset production code. The reusable 2D token renderer has
typed generational resources, managed texture/shader assets, painter-correct
primitives, and mock tests; the persistent \symbolref{SceneGraph} owns
generational hierarchy nodes and extracts borrowed render attachments through
that token path;
SparseCellGrid uses signed-coordinate 16$\times$16 chunks, retained candidate
and changed-frontier structures, and bounded worker evaluation for large
candidate/halo workloads; CanvasView uses CPU palette targets, \texttt{displayRgb}
fade, and one dirty-rect RGB display texture/quad; headless
\texttt{IllumoTests} and \texttt{IllumoGameTests} run through the aggregate
workspace build; the command console has measured editing,
metadata-driven help, and module-registered domain commands. Windows MSVC
Release/Debug builds are the supported target, DebugModule is excluded from
Release, and the Linux/macOS bootstraps are unsupported stale scaffolds.
\end{statusbox}

Debug console renderer diagnostics are \texttt{renderer\_demo [on|off]},
\texttt{assets}, and \texttt{asset\_reload <all|path>}.
